\documentclass[10pt]{article}
\usepackage{amsmath,amsfonts,amssymb,mathtools,mathrsfs,bm,bbm,array,xcolor,geometry}
\allowdisplaybreaks
\setlength{\parindent}{0pt}
\setlength{\parskip}{4pt}
\emergencystretch=4em

\geometry{paperwidth=215mm,paperheight=290mm,hmargin=15mm,vmargin=16mm}
\begin{document}
\section*{P0039 Drone Delivery Network Operations in a Circular Market}
\subsection*{Mathematical formulation}
\begin{align}
\max_{n,\omega,\beta}\quad
\Pi={}&2n\lambda mH_{1/2}(M)\sqrt{\frac{\beta\omega}{M}}
-2n\lambda\alpha\sqrt{\frac{\omega^5\beta^3}{M^3}}
\left[2H_{-1/2}(M)-H_{1/2}(M)\right]-kn^2,\tag{$P^M$}\\
\text{s.t.}\quad &n\beta\le\pi,\\
&n\ge0,\quad\omega\ge0,\quad\beta\ge0.
\end{align}

\subsection*{Definitions and model scope}
A retailer operates in a circular market of unit radius, with customers uniformly distributed around the circumference and symmetrically placed delivery centers. It chooses nonnegative center scale $n$, drone angular speed $\omega$ and service-arc length $\beta$ on each side of each center. Center scale is continuous and need not be an integer. Each side has a fixed positive integer number of delivery-time zones $M$. Demand coefficient $\lambda$ is strictly positive, while unit merchandise margin $m$, speed-squared coefficient $\alpha$ in distance delivery cost and network operating-cost coefficient $k$ are nonnegative.

Center scale multiplied by one-sided service length cannot exceed pi, preventing coverage beyond the circumference. Zero center scale means no operation and zero profit; it does not require division by zero centers. When scale, speed and coverage are all positive, each zone's promised time $\tau_i$ is its position number, starting at one, multiplied by one-sided service length and divided by the product of zone count and angular speed. At any positive promised time, demand per unit arc $D(\tau)$ is the demand coefficient divided by the square root of that promised time.

Operating profit $\Pi$ should be as high as possible. Over zone position numbers from first through last, calculate the sum of reciprocal square roots $H_{1/2}(M)$ and the sum of square roots $H_{-1/2}(M)$. Merchandise return is twice center scale times demand coefficient, unit margin and the reciprocal-square-root sum, multiplied by the square root of service length times speed divided by zone count. Delivery expense is twice center scale times demand coefficient and delivery-cost coefficient, multiplied by the square root of speed to the fifth power times service length cubed divided by zone count cubed, and then by twice the square-root sum minus the reciprocal-square-root sum. Subtract this expense and the network coefficient times squared center scale to obtain profit.

Apart from coverage and nonnegativity, scale, speed and service length have no added artificial bounds. There is no inventory ordering charge, finite-demand intercept or fixed speed charge. The positive-time demand interpretation applies only to positive operations; zero speed or zero coverage gives zero merchandise return through the stated products. If the supplied returns and costs make profit unbounded above, there is no finite best profit; an extra upper limit must not be invented.

\subsection*{Inline notation and relations}
$P^M$; $\bar P^M$; $M\in\mathbb Z_{++}$; $\lambda>0$; $\tau>0$; $D(\tau)=\lambda/\sqrt\tau$; $m\ge0$; $\alpha\ge0$; $k\ge0$; $H_r(M)=\sum_{i=1}^M i^{-r}$; $n\ge0$; $\omega\ge0$; $\beta\ge0$; $n$; $n>0$; $n\beta\le\pi$; $0\le\beta\le\pi/n$; $n=0$; $i$; $\tau_i=i\beta/(M\omega)$.
\end{document}